\chapter{Implementation}

The previous chapter described what the system should do and how its pieces fit
together. This chapter is about how those pieces were actually built: the setup
that lets the whole platform run on a laptop, the implementation choices that
turned the design into working code, and a few screenshots and generated samples
that show the result. The synthetic-data tool, MultiSeedGen, gets the most space,
because it is where most of the engineering effort went and it is the part of the
project we can defend in the most detail.

\section{Implementation environment and setup}

The whole system is containerized with Docker \cite{merkel2014docker}. The goal
was that a fresh checkout reaches a working stack with a single command, without
anyone installing a database or a Python toolchain on the host by hand. Docker
gives us an environment that behaves the same on a personal laptop, a lab
machine, or a server, which removes the usual "it works on my machine" problem
and makes it easy for a new team member to get started.

There are four primary components. The frontend is a React application \cite{reactindepth} that the user interacts with. The backend is a FastAPI service \cite{fastapi} that takes an image, runs the models, and returns the results. PostgreSQL stores the transactional records so a user can come back later and see the history of what they scanned. Alongside PostgreSQL, a ClickHouse instance acts as the analytics data warehouse, storing detailed inference telemetry and performance metrics for reporting and evaluation. The frontend talks to the backend over a simple HTTP API, and the backend communicates with the databases and object storage, so each part can be developed and tested on its own.

\autoref{fig:deployment} shows how the containers connect.

\figslot[1.1]{10-deployment.pdf}{Deployment view: the React frontend, the FastAPI backend, the PostgreSQL database, and the ClickHouse analytics warehouse.}{fig:deployment}

\section{Key implementation details}

======================================================

\subsection{Computer-vision background}

The computer vision pipeline is built to solve two distinct problems. The first is an inter-class problem: separating different seed types from one another within a mixed photograph. The second is an intra-class problem: addressing the quality differences between individual seeds of the same type. These correspond to object detection and per-seed classification, respectively. The following sections describe the model families used for each stage, leaving the detailed implementation and results to later chapters.

\subsection{Two-stage detectors}

Our detection model uses a ResNet backbone to extract strong visual features from the image, while the FPN and PANet help combine information from different feature levels so the model can detect both small and large objects better. In Faster R-CNN, the Region of Interest (ROI) part selects important areas from the image and classifies them more accurately. During training, CIoU is used to improve bounding box regression by making the predicted box overlap, center distance, and aspect ratio closer to the ground truth. Together, these components make the model more accurate and reliable for object detection.
ResNet backbone is the feature extractor at the start of the model. It takes the image and learns useful patterns like edges, shapes, and textures, so the detector has a strong representation of what is in the image.

FPN (Feature Pyramid Network) helps the model use features from different scales. This is important because small objects need fine details, while large objects need broader context, so FPN combines information from multiple levels of the network.

PANet (Path Aggregation Network) improves on FPN by strengthening the flow of information between low-level and high-level features. In simple terms, it helps the model pass details more effectively so detection becomes better, especially for objects of different sizes.

ROI (Region of Interest) refers to the important areas in the image that the detector focuses on after proposals are made. In Faster R-CNN, these regions are cropped and analyzed more carefully so the model can classify the object and refine its box.

CIoU (Complete Intersection over Union) is a loss function used to make bounding box predictions better. It does not just compare overlap, but also considers the distance between box centers and the box shape, which makes localization more accurate.

Faster R-CNN architecture is a two-stage object detection model. First, it finds possible object regions, then it classifies them and adjusts the boxes more precisely, which makes it accurate and reliable, though usually slower than one-stage models like YOLO.

\figslot[1.0]{fasterRCNN-architecture.pdf}{The Faster R-CNN detection architecture, showing the ResNet backbone, the FPN for multi-scale features, and the RPN and ROI head for detection and classification.}{fig:detector-arch}

The second stage is quality classification. To perform multi-label defect classification on individual isolated crops, the pipeline utilizes a deep convolutional neural network optimized for fine-grained feature extraction. Rather than relying on a standard off-the-shelf classifier, the architecture is a modified variant of EfficientNet-B2 integrated with a Convolutional Block Attention Module (CBAM) and a hybrid pooling layer.\cite{tan2019efficientnet} 

The system leverages an ImageNet-pretrained EfficientNet-B2 to serve as the primary feature extractor, CBAM Attention Layer is inserted immediately after the final feature map of the backbone. CBAM applies a sequential dual-attention mechanism. 

Channel Attention determines what defect characteristics to emphasize across channels, while Spatial Attention isolates where those defects are located on the physical seed geometry. Crucially, this attention block is wrapped in a residual connection to refine existing backbone features without degrading the network's underlying gradient stream.

Hybrid Pooling Head is the standard global average pooling layer is replaced with a concurrent, hybrid concatenation layout. By running Global Average Pooling (GAP) and Global Max Pooling (GMP) in parallel and splicing their outputs the network simultaneously tracks smooth structural global patterns alongside highly localized, sharp defect signals (such as microscopic fungal patches or fine cracks).

\figslot[1.0]{EfficientNet-architecture.pdf}{The EfficientNet-B2 defect classification architecture}{fig:detector-arch}

\subsection{One-stage detectors}

As a prominent class of one-stage architectures, single-shot detectors bypass the traditional region proposal phase entirely, enabling the simultaneous prediction of bounding box coordinates and corresponding class probabilities within a single forward pass. This unified approach was pioneered by YOLO \cite{redmon2016yolo}, which reformulated the object detection problem as a singular regression task mapped across a global grid framework. While this design strategy exchanges a marginal degree of spatial localization precision for a massive optimization in computational speed, it also yields strong generalization properties. Because the network processes images globally, it effectively learns robust representations that resist overfitting to localized environmental noise, thereby facilitating smoother transfer learning and superior domain adaptation when trained on synthetic datasets. 


\subsection{The two-stage detection and classification pipeline}

Checking a photo of seeds is really two separate problems, and the
implementation treats them as two models run one after the other. The first
stage is detection. An object detector scans the whole image and returns, for
every seed it finds, a bounding box, a confidence score, and the crop type, so the detector tells us where each seed is and which of the two it is. One photo therefore becomes a set of located, labelled seeds.

The detector is a Faster R-CNN \cite{fasterrcnn2015} with a ResNet-50 backbone \cite{resnet2015} and a Feature Pyramid Network \cite{lin2017fpn}. The feature pyramid lets the model find seeds at different sizes in the same image, which matters because a photo can hold a hundred seeds at slightly different distances from the camera. Faster R-CNN is the accurate option, while the second architecture, PANet, further strengthens feature flow across scales to improve detection of small and overlapping seeds. For situations where speed matters more than the last point of accuracy, such as a conveyor-belt feed or a phone in the field, there is an optional YOLO mode \cite{redmon2016yolo, yolov8} that runs the detection in a single pass and is much faster, at some cost in precision.

The second stage of the pipeline is fine-grained, multi-label classification. For each isolated seed, the system extracts the bounding box crop from the original image and routes it to a specialized quality network, which outputs predictions across a vector of concurrent defect conditions alongside their respective confidence scores. The core architecture leverages an ImageNet-pretrained EfficientNet-B2 backbone \cite{tan2019efficientnet}, selected for its optimal depth-width-resolution scaling balance and high computational throughput when processing dense seed batches.

\subsection{Model management and evaluation}

The backend also has a small amount of supporting machinery around the models.
It can register and swap the trained weights without code changes; it can run an
offline evaluation of a model against a held-out set, storing each run's metrics
in PostgreSQL alongside a Markdown report in object storage so that versions can
be compared before one is put into use; and it runs the actual inference in a
background worker so that uploading a batch of images does not block the rest of
the application while the models run. This part stayed simple on purpose, since
the prototype only needs to serve and compare a handful of models.

Choosing which model actually runs for a given request is the job of the model
resolver. Given a segment the kind of model needed (detection or classification) and, optionally, the seed type it resolves the model promoted
to \texttt{production} for that exact segment; if none has been promoted and a
seed type was supplied, it falls back to the global, seed-type-agnostic production
model for that kind; and if nothing is routable it reports that no model is ready
rather than failing silently. The decision is deterministic and stateless there
is no weighted routing and no per-user bucket so the same request always
resolves to the same model, and a developer can override the choice with an
explicit model on a single request. Figure~\ref{fig:model-resolution} shows the
decision.

\figslot[1.2]{16-model-resolution.pdf}{Model resolution: the resolver selects the model
promoted to production for a segment, falls back to the global production model
when only that is available, or reports that no model is ready.}{fig:model-resolution}

\subsection{Data warehouse and analytics telemetry (ClickHouse)}

To track the real-world performance of the grading system, the platform implements an online analytical processing (OLAP) star schema in ClickHouse. While PostgreSQL manages transactional and relational data (such as user credentials, scan batches, and model status), ClickHouse is dedicated to high-performance analytics, storing scan telemetry, inference latency, model accuracy over time, and geographic distribution metrics.

The analytics database schema consists of dimension tables (such as \texttt{dim\_user}, \texttt{dim\_model}, and \texttt{dim\_seed\_type}) and fact tables (such as \texttt{fact\_inference}, \texttt{fact\_detection}, and \texttt{fact\_scan\_batch}). Data is synchronized from PostgreSQL using application-level dual-writes rather than log-based change data capture (CDC). After each successful transactional commit, a Celery worker dispatches an idempotent insert to ClickHouse via the \texttt{AnalyticsRepository} using the async driver \texttt{clickhouse-connect}.

Because ClickHouse uses the \texttt{ReplacingMergeTree} table engine, duplicate inserts collapse on their sort key during merge cycles. This makes at-least-once delivery from the Celery tasks completely safe against duplicates, ensuring high consistency without complex distributed locking. Callers query this database on async read paths to drive the dashboard statistics, leaderboards, and offline performance tracking.

\subsection{MultiSeedGen: the synthetic-data pipeline}

Training a detector needs labelled images that contain many seeds at once, with a
box around each one. Those are far harder to collect than the single-seed photos
in the public datasets: someone would have to lay out dozens of seeds, photograph
them, and then hand-draw and label a box around every seed in every image.
MultiSeedGen avoids that work. It cuts individual seeds out of single-seed
source photos, pastes many of them onto realistic backgrounds, and exports a
fully labelled multi-seed detection dataset. Because the tool places every seed
itself, it already knows exactly where each seed is and what shape it has, so the
bounding box and the segmentation mask are known precisely. The annotations come
for free, with no manual labelling.

The hardest step is segmentation: cutting one seed cleanly out of its source
photo, including the soft edges. MultiSeedGen runs this as a cascade that starts
cheap and escalates only when needed. It first tries classical methods, a
colour-distance threshold against the background followed by GrabCut, which are
fast and good enough for seeds on a plain background. When those produce a poor
mask, it falls back to rembg, a U\textsuperscript{2}-Net model
\cite{qin2020u2net}, and for the hardest cases with feathered or low-contrast
edges it uses Segment Anything \cite{kirillov2023sam}. A cut-out whose mask
scores too low is dropped rather than pasted in with a bad outline, so a poor
segmentation never becomes a bad label.

With clean cut-outs in hand, the tool composites a scene. It pastes many seeds
onto a background, the cut-and-paste idea from Dwibedi et al.
\cite{dwibedi2017cutpaste}, varying the backgrounds so the detector does not just
learn one surface. To stop the synthetic images from looking too clean and too
similar to each other, MultiSeedGen applies domain-randomized augmentation: it
randomizes the geometry of each seed, its rotation, scale, and placement, and the lighting and colour of the scene \cite{tobin2017domainrand,
shorten2019augmentation, buslaev2020albumentations}. The aim is for a model
trained on these images to transfer to real photos, since it has seen a wide
range of arrangements and conditions rather than one narrow style.

Two details keep the resulting dataset trustworthy. First, the split between
training and validation is leakage-safe: a given source seed is assigned to
either training or validation, never both, so the same physical seed cannot show
up in a training image and a validation image at once. Without this, the
validation score would be inflated, because the model would effectively be tested
on seeds it had already seen. Second, the datasets export to both COCO and YOLO
formats \cite{lin2014coco}, the two formats the detection models expect, so the
generated data drops straight into training.

\section{Sample screenshots and output samples}

This section shows the running system and the results it produces.

The end product of the two-stage pipeline is an annotated image.
\autoref{fig:sample-output} shows the detector's boxes drawn back over a photo,
each one labelled with the classifier's good or bad verdict for that seed. This
is the overlay the frontend shows the user after a scan.

\figslot{sample-detection-output.png}{Annotated output: detection boxes with a
per-seed good or bad quality label from the classifier.}{fig:sample-output}

\autoref{fig:msg-sample} is a scene generated by MultiSeedGen, with the bounding
boxes drawn back from the exported labels. This is the kind of image the tool
produces by the thousand to train the detector.

\figslot{msg-sample-scene.png}{A MultiSeedGen scene: many cut-out seeds
composited onto a background, with bounding boxes drawn from the exported
labels.}{fig:msg-sample}

\autoref{fig:msg-seg-tuner} is the segmentation tuner from the MultiSeedGen
interface, which shows the kept and skipped cut-outs side by side so we can see
which seeds a given method dropped, and why, before running a full generation.

\figslot{msg-seg-tuner.png}{MultiSeedGen segmentation tuner: kept and skipped
cut-outs shown together so the segmentation quality can be checked before a full
run.}{fig:msg-seg-tuner}


------------------------

The end product of processing an entire collection sample is the batch classification summary. \autoref{fig:batch-output-summary} illustrates the structured batch report generated after evaluating dense arrays of seeds, tracking unique filenames alongside their total detected element count. This output summarizes the core execution loop by coordinating individual localized crops into a single database format.

\figslot{batch-output-summary.png}{Batch processing summary output: an evaluation overview mapping scanned sample images to their respective total detected seed counts.}{fig:batch-output-summary}

{fig:seed-defect-analytics} displays the analytical dashboard tracking the total percentage of bad seeds and the ratio of individual defect categories across all processed images.

\figslot{seed-defect-analytics.png}{Statistical distribution analytics: analytical graphs illustrating the total percentage of bad seeds and detailed defect categories scanned across the entire batch.}{fig:seed-defect-analytics}

To handle multi-crop families seamlessly, the infrastructure maintains distinct directory paths for optimized weights. \autoref{fig:saved-models-directory} outlines the serialized model store housing active checkpoints for both the object detection layers and the fine-grained multi-label classification backbones. Detections are dynamically grouped by crop type and dispatched to these corresponding weights.

\figslot{saved-models-directory.png}{Saved models repository: the file directory tree showing isolated checkpoint weights grouped by respective crop families and execution stages.}{fig:saved-models-directory}

At a granular level, validation accuracy depends entirely on distinct feature boundaries. \autoref{fig:good-bad-crops} shows a curated compilation of the extracted seed crops, showcasing the structural contrast between healthy seeds and degraded samples. This highlights how the model decouples natural biological variance from physical deformities.

\figslot{good-bad-crops-comparison.png}{Granular sample classification: isolated good seeds and bad seeds side by side, showcasing the clear morphological differences detected by the model.}{fig:good-bad-crops}